\documentclass[10pt,a4paper]{amsart}
\usepackage[margin=19mm]{geometry}
\usepackage{amsmath,amssymb,mathtools}
\usepackage[scr=rsfs,cal=euler]{mathalfa}
\usepackage{xfrac}
\usepackage[unicode,hidelinks,bookmarks=false]{hyperref}
\newcommand{\ie}{i.e.\ }
\newcommand{\cf}{cf.\ }
\newcommand{\resp}{respectively\ }
\newcommand{\Subsection}{\subsection}
\providecommand{\nobreakdash}{\nobreak\hbox{-}}
\setlength{\emergencystretch}{2em}
\multlinegap0pt
\usepackage{amssymb}
\newtheorem{definition}{Definition}
\newtheorem{lemma}{Lemma}
\newtheorem{proposition}{Proposition}
\newtheorem{corollary}{Corollary}
\newtheorem{theorem}{Theorem}
\title{On the uniqueness and the regularity of solutions to the Navier Stokes in the Lorentz space $L^{3,\infty}(\mathbb{R}^3)$}
\author{Ramzi May}
\date{Source: arXiv:2609.14688v1 (CC BY 4.0)}
\begin{document}
\maketitle

\noindent\emph{Source and licence.} On the uniqueness and the regularity of solutions to the Navier Stokes in the Lorentz space $L^{3,\infty}(\mathbb{R}^3)$. Version arXiv:2609.14688v1 (CC BY 4.0), distributed by arXiv under the Creative Commons Attribution 4.0 International licence, http://creativecommons.org/licenses/by/4.0/, as recorded in the arXiv licence metadata for this exact version and shown on the abstract page https://arxiv.org/abs/2609.14688v1. Full licence notice: \texttt{licenses/arxiv-2609.14688-CC-BY-4.0.txt}.\par\smallskip

\noindent\textit{Editorial scope.} Counted unit: the article's theorem WS (source0.tex lines 621--623). The article's complete original proof of it is delivered with it (source0.tex lines 624--810, one \texttt{proof} environment of 187 lines). No mathematics is written, repaired or re-derived: the statement and the proof are the article's own lines, in the article's own notation, and no retained line is substituted or re-flowed.  Retained uncounted as the source-local context the proof needs: lines 152--164; lines 257--278; lines 324--330; lines 332--334; lines 343--390; lines 412--418; lines 429--438; lines 456--460; lines 465--469; lines 521--523; lines 548--552; lines 578--586.  Nothing was omitted from any retained span, so the package carries no omission record.  No retained line contains a citation, so the wrapper carries no bibliography and no bibliography entry.  No retained line contains a figure environment, an image-inclusion command or drawing code, so no image is delivered and the package declares no asset dependency.  Editorial formatting only, in addition: an AMS \texttt{amsart} preamble that restates the article's own declarations and macros used by the retained lines, namely the environments \texttt{definition}, \texttt{lemma}, \texttt{proposition}, \texttt{corollary}, \texttt{theorem} on separate counters, together with the standard packages those lines need.  The retained environments print in this document as Definition 1, Lemma 1, Proposition 1, Corollary 1, Theorem 1 in order of appearance, whereas the article prints the same items with the numbering of its own full document.  The complete native source of the article as distributed in the arXiv e-print for this exact version is bundled unchanged as \texttt{sources/210348/source0.tex}.
\begin{definition}\label{reg} [admissible spaces]
A Banach space $X$ is called admissible space if it satisfies the following three
properties:
\begin{enumerate}
\item $\mathcal{S}(\mathbb{R}^{n})\hookrightarrow X\hookrightarrow \mathcal{S}^\prime(\mathbb{R}^{n})$.
\item $X\subset L^1_{\text{loc}}(%
\mathbb{R}^{n}).$
\item For every $f\in X$ and $x_0\in \mathbb{R}^{n}$,$f(.+x_0)\in X$ and $\left\Vert f(.+x_0)\right\Vert _{X}=\left\Vert  f\right\Vert _{X}.$
\item For every $f\in X$ and $h\in L^\infty(\mathbb{R}^{n})$, $hf\in X$ and $\left\Vert hf\right\Vert _{X}\lesssim\left\Vert h\right\Vert
_{\infty }\left\Vert f\right\Vert _{X}$.
\end{enumerate}
We denote by $\tilde{X}$ the  closure of $\mathcal{S}(\mathbb{R}^{n})$ in $X$.
\end{definition}

\begin{lemma}\label{para}
Let $X$ be an admissible space. Then
\begin{enumerate}
\item For any $s>0,$ the bilinear operator $\pi ^{b}$ is continuous from $%
L^{\infty }(\mathbb{R}^{n})\times B_{X}^{s,\infty }(\mathbb{R}^{n})$
into $B_{X}^{s,\infty }(\mathbb{R}^{n}).$

\item For any real number $s,$ the bilinear operator $\pi ^{a}$ is continuous
from $L^{\infty }(\mathbb{R}^{n})\times B_{X}^{s,\infty }(\mathbb{R}^{n})
$ into $B_{X}^{s,\infty }(\mathbb{R}^{n}).$

\item For any real number $s_{1}<0$ and $s_{2}$ such that $%
s_{1}+s_{2}>0$, the bilinear operator $\pi ^{b}$ is continuous from $%
B_{\infty }^{s_{1},\infty }(\mathbb{R}^{n})\times B_{X}^{s_{2},\infty }(%
\mathbb{R}^{n})$ into $B_{X}^{s_{1}+s_{2},\infty }(\mathbb{R}^{n}).$

\item For any negative real number $s_{1}$ and any real number $s_{2}$, the
bilinear operator $\pi ^{a}$ is continuous from $B_{\infty }^{s_{1},\infty }(%
\mathbb{R}^{n})\times B_{X}^{s_{2},\infty }(\mathbb{R}^{n})$ into $%
B_{X}^{s_{1}+s_{2},\infty }(\mathbb{R}^{n}).$
\end{enumerate}
\end{lemma}

\begin{proposition} [Maximal $L^p(L^q)$ regularity for the heat kernel] \label{max}
Let $T>0 $ and $1<p,q<\infty$. Then
\[
\left\Vert \Delta \int_{0}^{t}e^{(t-\tau )\Delta }f(\tau )d\tau \right\Vert_{L^p([0,T],L^q(\mathbb{R}^n))}\lesssim\left\Vert f \right\Vert_{L^p([0,T],L^q(\mathbb{R}^n))}
\]
for every $f\in L^p([0,T],L^q(\mathbb{R}^n))$.
\end{proposition}

\begin{lemma}\label{Fou}
Let $X$ be an admissible space, $s$ a real number and $1\leq q\leq\infty$. If $\Theta:\mathbb{R}^n_\ast\to\mathbb{R}$ is a $C^\infty$  homogenous function of degree $\alpha>0$ (i.e., $\Theta(\lambda\xi)=\lambda^\alpha\Theta(\xi), \forall \lambda>0, \xi\neq 0$) then the Fourier multiplier $\Theta(D)$ is continuous from $B^{s,q}_X$ to $B^{s-\alpha,q}_X$. In particular, the operator $\mathbb{P}\nabla$ is continuous from $B^{s,q}_X$ to $B^{s-1,q}_X$.
\end{lemma}

\begin{proposition}\label{lorentz} [Some properties of Lorentz spaces]
Let $1<p<\infty $ and $1\leq q\leq \infty .$ The following assertions hold true.
\begin{enumerate}
\item $L^{p,q}(\mathbb{R}^{n})$ is an admissible space.
\item $v\in L^{p,\infty }(\mathbb{R}^{n})$ if and only if $\sup_{\lambda
>0}\lambda ^{p}m(\{x\in \mathbb{R}^{n}:\left\vert v(x)\right\vert \geq
\lambda \})<\infty $ where $m$ is the Lebesgue measure on $\mathbb{R}^{n}$. Moreover
\[
\left\Vert v\right\Vert_{L^{p,\infty }(\mathbb{R}^{n})}\simeq \sup_{\lambda
>0}\lambda [m(\{x\in \mathbb{R}^{n}:\left\vert v(x)\right\vert \geq
\lambda \}]^{1/p}.
\]
\item Let $0<\alpha <n$. n the  function $f(x)=\frac{1}{\left\vert x\right\vert ^{\alpha }%
}$ belongs to the space $L^{\frac{n}{\alpha },\infty }(%
\mathbb{R}^{n}).$

\item $L^{p,1}(\mathbb{R}^{n})\hookrightarrow L^{p,q}(\mathbb{R}%
^{n})\hookrightarrow L^{p,\infty }(\mathbb{R}^{n}).$

\item $L^{p,p}(\mathbb{R}^{n})=L^{p}(\mathbb{R}^{n}).$
\item If $1<p_1<p<p_2<\infty$ then $L^{p_1,\infty}(\mathbb{R}^n)\cap L^{p_2,\infty}(\mathbb{R}^n) \hookrightarrow L^p(\mathbb{R}^{n}).$

\item $\left\Vert f(\lambda .)\right\Vert _{L^{p,q}(\mathbb{R}^{n})}=\lambda
^{-\frac{n}{p}}\left\Vert f\right\Vert _{L^{p,q}(\mathbb{R}^{n})}~\forall
\lambda >0~\forall f\in L^{p,q}(\mathbb{R}^{n}).$

\item $L^{p,\infty }(\mathbb{R}^{n})\hookrightarrow B_{\infty }^{-\frac{n}{p}%
,\infty }(\mathbb{R}^{n}).$

\item If $1<p_{1},p_{2},p_{3}<\infty $ and $1\leq q_{1},q_{2},q_{3}\leq
\infty $ such that $\frac{1}{p_{3}}=\frac{1}{p_{1}}+\frac{1}{p_{2}}$ and $%
\frac{1}{q_{3}}=\frac{1}{q_{1}}+\frac{1}{q_{2}}$ then%
\[
\left\Vert fg\right\Vert _{L^{p_{3},q_{3}}}\leq \left\Vert f\right\Vert
_{L^{p_{1},q_{1}}}\left\Vert g\right\Vert _{L^{p_{2},q_{2}}}\forall f\in
L^{p_{1},q_{1}}(\mathbb{R}^{n}),g\in L^{p_{2},q_{2}}(\mathbb{R}^{n}).
\]

\item If $1<p_{1},p_{2},p_{3}<\infty $ and $1\leq q_{1},q_{2},q_{3}\leq
\infty $ such that $\frac{1}{p_{3}}=\frac{1}{p_{1}}+\frac{1}{p_{2}}-1$ and $%
\frac{1}{q_{3}}=\frac{1}{q_{1}}+\frac{1}{q_{2}}$ then%
\[
\left\Vert f\ast g\right\Vert _{L^{p_{3},q_{3}}}\leq \left\Vert f\right\Vert
_{L^{p_{1},q_{1}}}\left\Vert g\right\Vert _{L^{p_{2},q_{2}}}\forall f\in
L^{p_{1},q_{1}}(\mathbb{R}^{n}),g\in L^{p_{2},q_{2}}(\mathbb{R}^{n}).
\]
\end{enumerate}
\end{proposition}

\begin{corollary}\label{Mon}
Let $f\in L^{n,\infty }(\mathbb{R}^{n})$ and $g\in L^{q}(\mathbb{%
R}^{n})$  with $\frac{n}{n-1}<q<\infty$. Then%
\[
\left\Vert \frac{1}{\sqrt{-\Delta }}(fg)\right\Vert _{q}\lesssim \|f\|_{L^{n,\infty}} \|g\|_{q}.
\]
\end{corollary}

\begin{corollary} \label{cor2}
Let $T>0$ and $ 1<\rho _{1}<2.$ Then for every $f\in L^{\rho _{1}}([0,T[;%
\mathbb{R}),$
\[
\left\Vert \int_{0}^{t}\frac{f(s)}{\sqrt{t-s}}ds\right\Vert _{L^{\rho
_{2}}([0,T[,\mathbb{R})}\lesssim \left\Vert f\right\Vert _{L^{\rho _{1}}([0,T[,%
\mathbb{R})}
\]%
where $\rho _{2}=\frac{2\rho _{1}}{2-\rho _{1}}$.
\end{corollary}

\begin{lemma}\label{con}
Let $T>0$ and $1\leq p\leq\infty$. For every $f\in L^{p}([0,T[;%
\mathbb{R}),$
$$ \left\Vert T(f)\right\Vert_{L^{p}([0,T[,\mathbb{R})}\leq 2\sqrt{T}\left\Vert f\right\Vert_{L^{p}([0,T[,\mathbb{R})}.$$
\end{lemma}

\begin{lemma}\label{GMO}
Let $1<p<\infty$. For any $f\in W^{1,p}(\mathbb{R}^n)\cap B^{-1,\infty}_\infty(\mathbb{R}^n)$,
$$\left\Vert f\right\Vert_{2p}\lesssim \left\Vert f\right\Vert_{W^{1,p}}^{\frac{1}{2}}\left\Vert f\right\Vert_{B^{-1,\infty}_\infty }^{\frac{1}{2}} $$
where $\left\Vert f\right\Vert_{W^{1,p}}=\left\Vert f\right\Vert_p +\left\Vert\sqrt{-\Delta} f\right\Vert_p$.
\end{lemma}

        \begin{equation}\label{kernel}
            \left\Vert \mathcal{K}_t\right\Vert_{ L^1(\mathbb{R}^3)}=\frac{c}{\sqrt{t}},
        \end{equation}

\begin{lemma}\label{Linj}
If $X$ is a Kato space then $X\hookrightarrow B_{\infty }^{-1,\infty }(%
\mathbb{R}^{3}).$ More precisely $B^{0,\infty}_X\hookrightarrow B_{\infty }^{-1,\infty }(%
\mathbb{R}^{3}).$
\end{lemma}

\begin{proposition}\label{Ram}
Let $X$ be a Kato space and let $u_0\in\tilde{X}$. there exists a unique maximal solution $u\in C([0,T^\ast),\tilde{X})\cap C_{0,\frac{1}{2}}([0,T^\ast),L^\infty(\mathbb{R}^n))$ to the Navier Stokes equations with initial data $u_0$ where $0<T^\ast\leq \infty$ is the maximal time of existence of the solution $u$. Moreover the following proprieties are satisfied:
\begin{enumerate}
    \item For every $s>0$, $u\in C^\infty((0,T^\ast),B^{s,\infty}_X\cap B^{s,\infty}_\infty).$
    \item It $T^\ast< \infty$ then $u(t)$ does not converge in $B^{-1,\infty}_\infty(\mathbb{R}^n)$ as $t\to T^{\ast -}$.
    \item For every $\alpha>0, u\in C_{0,\frac{\alpha}{2}}([0,T^\ast),B^{\alpha,\infty}_X)$.
\end{enumerate}
Such solution $u$ will be called The $X-$Kato maximal solution to the Navier-Stokes equations associated to the initial data $u_0$.
\end{proposition}

\begin{theorem}\label{WS}
Let $v_0\in\tilde{L}^{3,\infty}(\mathbb{R}^3)$, $v\in C([0,T^\ast),\tilde{L}^{3,\infty}(\mathbb{R}^3))\cap C_{0,\frac{1}{2}}([0,T^\ast),L^\infty(\mathbb{R}^3))$ be the  $ {L}^{3,\infty}(\mathbb{R}^3)$-Kato maximal solution to the Navier Stokes equations with initial data $v_0$, and $u\in L^p([0,T_1],\tilde{L}^{3,\infty}(\mathbb{R}^3))\cap C([0,T_1],B^{-1,\infty}_\infty(\mathbb{R}^3))$, with $p>2$ and $0<T_1<\infty$, a solution to the Navier Stokes equations associated to the same initial data $v_0$. Then $T^\ast>T_1$ and $u(t)=v(t)$ on $[0,T_1]$.
\end{theorem}

\begin{proof}
Let $0<T<\min{\{T_1,T^\ast\}}.$ First, we may assume, without loss of generality, that
$p\in(2,4)$. Set $w_{1}=B(u,u),w_{2}=B(v,v)$, and $w=w_{1}-w_{2}.$ We shall
prove that $w=0$ on $[0,T]$ which implies $u=v$ on $[0.T]$ since $w$ is also equal to   $u-v$. The proof is divided  into several steps:
\par\noindent\textbf{Step 1:} We shall prove that there exists $q>2$
such that $w_{1},w_{2}\in L^{p}([0,T],L^{q}(\mathbb {R}^{3})).$
\par\noindent Using the estimation (\ref{kernel}) on the Oseen's kernel, the product estimate in Lorentz spaces (Proposition \ref{lorentz}), and the fact that Lorentz spaces are admissible spaces,  we obtain
\[
\left\Vert w_{1}(t)\right\Vert _{L^{\frac{3}{2},\infty }(\mathbb{R}%
^{3})}\lesssim \int_{0}^{t}\frac{1}{\sqrt{t-s}}\left\Vert u(s)\right\Vert
_{L^{3,\infty }(\mathbb{R}^{3})}^{2}ds.
\]%
Therefore, in view of  Corollary \ref{cor2}, we have
\[
\left\Vert w_{1}\right\Vert _{L^{\rho }([0,T],L^{\frac{3}{2},\infty }(%
\mathbb{R}^{3}))}\lesssim \left\Vert u\right\Vert _{L^{p}([0,T],L^{3,\infty
}(\mathbb{R}^{3}))}^{2}
\]%
where $\rho =\frac{2p}{4-p}.$ This implies that $w_1\in L^{p }([0,T],L^{\frac{3}{2},\infty }( \mathbb{R}^{3}))$ since $p\leq \rho$. On the other hand, since $%
e^{t\Delta }v_{0}\in C([0,T],\tilde{L}^{3,\infty }(\mathbb{R}^{3})),$ $%
w_{1}=u-e^{t\Delta }v_{0}\in L^{p}([0,T],\tilde{L}^{3,\infty }(\mathbb{R}%
^{3}))$, then using the continuous embedding of the space  $L^{3,\infty }(\mathbb{R}^{3})\cap L^{%
\frac{3}{2},\infty }(\mathbb{R}^{3})$ into $
L^{q}(\mathbb{R}^{3})$ for any $q\in (\frac{3}{2},3),$  we conclude that $%
w_{1}\in L^{p}([0,T],L^{q}(\mathbb{R}^{3}))$ for any real number $q\in (\frac{3}{2},3).$ The same result
holds true for $w_{2}$, since $v\in C([0,T],\tilde{L}^{3,\infty }(\mathbb{R}%
^{3})\subset L^{p}([0,T],\tilde{L}^{3,\infty }(\mathbb{R}^{3})).$

\par\noindent\textbf{Step 2:} We shall  show that
\begin{equation} \label{claim1}
\lim_{t\to 0^+} B(w,w)(t)=0  \text{ in } B_{\infty }^{-1,\infty }.
\end{equation}
Since
\begin{equation}\label{aicha}
w=u-v\to 0 \text{ in } B_{\infty }^{-1,\infty } \text{ as } t\to 0^+,
\end{equation}
and $w=B(w,w)+2B(v,w)$, it suffices to prove that
\begin{equation}\label{claim}
B(v,w)(t)\to 0  \text{ in } B_{\infty }^{-1,\infty } \text{ as } t\to 0^+   .
\end{equation}
To do this, we use  the Bony-paraproduct operators to split $B(v,w)(t)$ into two parts:
\begin{eqnarray*}
B(v,w)(t) &=&\int_{0}^{t}\mathbb{P}\nabla e^{(t-s)\Delta }\pi
^{a}(v,w)(s)ds+\int_{0}^{t}\mathbb{P}\nabla e^{(t-s)\Delta }\pi ^{b}(w,v)(s)
\\
&=&W_a(t)+W_{b}(t).
\end{eqnarray*}%
Using now the fourth assertion of Lemma \ref{para} and the last assertion of Lemma \ref{reg}, we get
\[
\left\Vert W_a(t)\right\Vert _{B_{\infty }^{-1,\infty }}\lesssim
\sup_{0\leq s\leq t}\left\Vert v(s)\right\Vert _{B_{\infty }^{-1,\infty
}}\sup_{0\leq s\leq t}\left\Vert w(s)\right\Vert _{B_{\infty }^{-1,\infty
}},
\]%
which, combined with (\ref{aicha}),  implies that
\begin{equation}\label{claim2}
\left\Vert W_a(t)\right\Vert _{B_{\infty }^{-1,\infty }}\rightarrow 0\text{
as }t\rightarrow 0.
\end{equation}
Let $1<\alpha <2$ be a fixed real number. Using  Lemma \ref{reg}, Lemma \ref{para}, and Lemma \ref{Fou}, we obtain
\begin{eqnarray*}
\left\Vert W_{b}(t)\right\Vert _{B_{L^{3,\infty }}^{0,\infty }} &\lesssim
&\int_{0}^{t}\frac{1}{(t-s)^{1-\frac{\alpha }{2}}}\left\Vert \mathbb{P}%
\nabla \pi ^{b}(w,v)(s)\right\Vert _{B_{L^{3,\infty }}^{\alpha -2,\infty }}ds
\\
&\lesssim &\int_{0}^{t}\frac{1}{(t-s)^{1-\frac{\alpha }{2}}}\left\Vert \pi ^{b}(w,v)(s)\right\Vert _{B_{L^{3,\infty }}^{\alpha -1,\infty }}ds \\
&\lesssim &\int_{0}^{t}\frac{1}{(t-s)^{1-\frac{\alpha }{2}}}\left\Vert
w(s)\right\Vert _{B_{\infty }^{-1,\infty }}\left\Vert
v(s)\right\Vert _{B_{L^{3,\infty }}^{\alpha ,\infty }}ds \\
&\lesssim &\int_{0}^{t}\frac{1}{(t-s)^{1-\frac{\alpha }{2}}}\frac{1}{s^{%
\frac{\alpha }{2}}}ds\sup_{0\leq s\leq t}\left\Vert
w(s)\right\Vert _{B_{\infty }^{-1,\infty }} \sup_{0\leq s\leq t}s^{\frac{\alpha }{2}}\left\Vert v(s)\right\Vert
_{B_{L^{3,\infty }}^{\alpha ,\infty }}\\
&=&C(\alpha ) \sup_{0\leq s\leq t}\left\Vert w(s)\right\Vert
_{B_{\infty }^{-1,\infty }}\sup_{0\leq s\leq t}s^{\frac{\alpha }{2}}\left\Vert v(s)\right\Vert _{B_{L^{3,\infty
}}^{\alpha ,\infty }}.
\end{eqnarray*}%
Combining the last inequality with the continuous embedding $B_{L^{3,\infty }}^{0,\infty }\hookrightarrow
B_{\infty }^{-1,\infty }$ (see Lemma \ref{Linj}),  we infer that $\left\Vert W_a(t)\right\Vert_{B_{\infty }^{-1,\infty }}\to 0$ as $t\to o^+$ . This combined with (\ref{claim2}) finishes the proof of the required claim (\ref{claim}).
\par\noindent\textbf{Step 3:}
Let $\tau \in (0,T]$, First, by the estimate (\ref{kernel}),
\[
\left\Vert B(w,w)(t)\right\Vert_{
L^{\frac{q}{2}}(\mathbb{R}^{3})}\lesssim \int_{0}^{t}\frac{1}{%
\sqrt{t-s}}\left\Vert w(s)\right\Vert_{L^{q}(\mathbb{R}^{3})}^2ds,
\]%
which implies, thanks to Corollary \ref{cor2},
\[
\left\Vert B(w,w)\right\Vert _{L^\rho([0,\tau ],L^{\frac{q}{2}}(%
\mathbb{R}^{3}))}\lesssim \left\Vert w\right\Vert
_{L^{p}([0,\tau ],L^{q}(\mathbb{R}^{3}))}^{2},
\]%
where $\rho =\frac{2p}{4-p}.$ Since $p\leq\rho$,
the previous estimate gives
\begin{equation}\label{he1}
\left\Vert B(w,w)\right\Vert _{L^p([0,\tau ],L^{\frac{q}{2}}(%
\mathbb{R}^{3}))}\lesssim \left\Vert w\right\Vert
_{L^{p}([0,\tau ],L^{q}(\mathbb{R}^{3}))}^{2}.
\end{equation}
On the other hand,
\[
\sqrt{-\Delta } B(w,w)(t)=\frac{\sqrt{-\Delta }\mathbb{P}\nabla }{\Delta } \Delta \int_{0}^{t} e^{(t-s)\Delta }(w w)(s)ds.
\]
Therefore, From Proposition \ref{max} and the continuity of the operator $%
\frac{\sqrt{-\Delta }\mathbb{P}\nabla }{\Delta }$ on the Lebesgue space $L^{%
\frac{q}{2}}(\mathbb{R}^{3}),$ we deduce that
\begin{equation}\label{he2}
\left\Vert \sqrt{-\Delta }B(w,w)\right\Vert _{L^{\frac{p}{2}}([0,\tau ],L^{%
\frac{q}{2}}(\mathbb{R}^{3}))}\lesssim \left\Vert w\right\Vert
_{L^{p}([0,\tau ],L^{q}(\mathbb{R}^{3}))}^{2}.
\end{equation}
Combining the estimates (\ref{he1}) and (\ref{he2}), we get
\[
\left\Vert B(w,w)\right\Vert _{L^{\frac{p}{2}}([0,\tau ],W^{1,\frac{q}{2}}(%
\mathbb{R}^{3}))}\lesssim \left\Vert w\right\Vert _{L^{p}([0,\tau ],L^{q}(%
\mathbb{R}^{3}))}^{2}.
\]%
Hence, applying Lemma \ref{GMO}, we obtain
\begin{equation}
\left\Vert B(w,w)\right\Vert _{L^{p}([0,\tau ],L^{q}(\mathbb{R}%
^{3}))}\lesssim \left(\sup_{0\leq s\leq \tau }\left\Vert B(w,w)(s)\right\Vert
_{B_{\infty }^{-1,\infty }}\right)^\frac{1}{2}\left\Vert w\right\Vert _{L^{p}([0,\tau ],L^{q}(%
\mathbb{R}^{3}))}.  \label{Est3}
\end{equation}%
Let $\varepsilon >0$ to be fixed later. The density of $L^{\infty }(%
\mathbb{R}^{3})\cap L^{3,\infty }(\mathbb{R}^{3})$ in $\tilde{L}^{3,\infty }(%
\mathbb{R}^{3})$ and the fact that $v\in C([0,T],\tilde{L}^{3,\infty }(%
\mathbb{R}^{3}))$ ensure the existence of $v_{\varepsilon }\in
C([0,T],L^{3,\infty }(\mathbb{R}^{3}))$ and $~v_{\infty }\in
C([0,T],L^{\infty }(\mathbb{R}^{3}))$ such that
\[
v=v_{\varepsilon }+v_{\infty },\sup_{0\leq s\leq T}\left\Vert v_{\varepsilon
}(s)\right\Vert _{L^{3,\infty }}\leq \varepsilon \text{ and }\sup_{0\leq
s\leq T}\left\Vert v_{\varepsilon }(s)\right\Vert _{\infty }=M_{\infty
}<\infty .
\]%
We decompose $B(v,w)$ into two terms as follows
\begin{eqnarray*}
B(v,w)(t)&=&
\int_{0}^{t} \mathbb{P}\nabla e^{(t-s)\Delta }(v_{\varepsilon
}w)(s)ds+\int_{0}^{t}\mathbb{P}\nabla e^{(t-s)\Delta }(v_{\infty
}w)(s)ds\\
&=& \frac{\mathbb{P}\nabla\sqrt{-\Delta }}{\Delta }
\Delta \int_{0}^{t}e^{(t-s)\Delta }\frac{1}{\sqrt{-\Delta }}(v_{\varepsilon
}w)(s)ds+\int_{0}^{t}\mathbb{P}\nabla e^{(t-s)\Delta }(v_{\infty
}w)(s)ds\\
&=& \Sigma_{\varepsilon}(t)+\Sigma_{\infty }(t).
\end{eqnarray*}
By applying, respectively,  the continuity of the operator $\frac{\mathbb{P}\nabla\sqrt{-\Delta }}{\Delta}$ on $ L^{q}(\mathbb{R}^{3})$, Proposition \ref{max} and Corollary \ref{Mon},  we obtain
\begin{eqnarray}
\left\Vert \Sigma_{\varepsilon}\right\Vert _{L^{p}([0,\tau ],L^{q}(\mathbb{R}%
^{3}))} &\lesssim &\sup_{0\leq s\leq \tau }\left\Vert v_{\varepsilon
}(s)\right\Vert _{L^{3,\infty }}\left\Vert w\right\Vert _{L^{p}([0,\tau
],L^{q}(\mathbb{R}^{3}))}  \nonumber \\
&\lesssim &\varepsilon \left\Vert w\right\Vert _{L^{p}([0,\tau ],L^{q}(%
\mathbb{R}^{3}))}.  \label{Est2}
\end{eqnarray}%
On the other hand, using as usual the estimate (\ref{kernel}) on the Ossen kernel, we get
\[\left\Vert \Sigma_{\infty }(t)\right\Vert _{L^{q}(\mathbb{R}^{3})} \lesssim
\int_{0}^{t}\frac{1}{\sqrt{t-s}}M_{\infty }\left\Vert w(s)\right\Vert
_{L^{q}(\mathbb{R}^{3})}ds
\]
Therefore, from Lemma \ref{con}, we have
\begin{equation}\label{Est0}
\left\Vert \Sigma_{\infty }\right\Vert _{L^{p}([0,\tau ],L^{q}(\mathbb{R}^{3}))}\lesssim M_{\infty }\sqrt{\tau }\left\Vert w\right\Vert_{L^{p}([0,\tau ],L^{q}(\mathbb{R}^{3}))}
\end{equation}
Finally, going back to the identity $w=B(w,w)+2B(v,w)$ and gathering the
estimates (\ref{Est3}), (\ref{Est2}), and (\ref{Est0}), we infer that
there exists an absolute constant $C>0$ (independent of $\varepsilon $ and $%
\tau $) such that for any $\tau \in \lbrack 0,T]$%
\[
\left\Vert w\right\Vert _{L^{p}([0,\tau ],L^{q}(\mathbb{R}^{3}))}\leq
C\big((\sup_{0\leq s\leq \tau }\left\Vert B(w,w)(s)\right\Vert
_{B_{\infty }^{-1,\infty }})^\frac{1}{2}+\varepsilon +M_{\infty }\sqrt{\tau }\big) \left\Vert
w\right\Vert _{L^{p}([0,\tau ],L^{q}(\mathbb{R}^{3}))}.
\]%
Hence, using the estimate (\ref{claim1}) and  choosing respectively $\varepsilon $ and $\tau $ small enough, we conclude that there exists $\tau \in (0,T]$ such
that
\[
\left\Vert w\right\Vert _{L^{p}([0,\tau ],L^{q}(\mathbb{R}^{3}))}\leq \frac{1%
}{2}\left\Vert w\right\Vert _{L^{p}([0,\tau ],L^{q}(\mathbb{R}^{3}))}
\]%
which implies that $w(t)=0$ on $[0,\tau ]$ and by consequent $u(t)=v(t)$ on $%
[0,\tau ].$ We shall now apply a well-known iteration argument to infer that  $u(t)=v(t)$ on the entire interval $[0,T].$ Let $t_\ast=\sup\{0\leq t\leq T: u=v  \text{ on } [0,t]\}$. We have $t_\ast\in[\tau,T].$ If $t_\ast<T$, then $u(.+t_\ast)$ and $v(.+t_\ast)$ are two mild solutions on $[0,T-t_\ast]$ to the Navier-Stokes equations associated to the same initial data $u(t_\ast)=v(t_\ast)$. Hence, applying what we have already done for $u$ and $v$ ensures the existence of $\tau_\ast\in(0,T-t_\ast)$ such that $u(.+t_\ast)=v(.+t_\ast)$ on the interval $[0,\tau_\ast ]$. This contradicts the definition of $t_\ast$. We then conclude that $u(t)=v(t)$ on $[0,T]$. Since the argument holds for every  $ T\in (0,\min\{T^\ast,T_1\}),$ we infer  that $u(t)=v(t)$ on the greater interval $[0,\min\{T^\ast,T_1\}).$ Finally, necessary  $T^\ast> T_1$, because if $T^\ast\leq T_1$ then
$$\lim_{t\to T^{\ast -}}v(t)=\lim_{t\to T^{\ast -}}u(t)=u(T^\ast)$$
in $B^{-1,\infty}_\infty(\mathbb{R}^3)$, which is impossible in view of  the last assertion of Proposition \ref{Ram}. This completes  the proof.
\end{proof}

\end{document}
